\subsection{The integral category and finite comparison}

The expression ``integral Betti fibre functor'' requires an integral source. On a nonzero $\Q$-linear category, multiplication by $n$ is an invertible morphism for every positive integer $n$. Its image under an additive functor to free $\Z$-modules would be multiplication by $n$, which is not invertible for $n>1$. Thus no such functor can literally have that source and target. In this section our integral comparisons concern the concrete diagram of smooth projective varieties, integral cohomology modulo torsion, and integral algebraic correspondences, or any integral motivic realization extending it. The counterexample below already occurs on abelian varieties with actual homomorphisms and covariant $H_1(-,\Z)$. Dualizing free lattices gives the identical obstruction for contravariant $H^1$. This states the precise restriction map from a proposed motivic comparison to the diagram we test.

\begin{result}{E.1.1}{Finite diagrams and finitely generated fields}{Proved}
Let $k\subset\C$ have characteristic zero, and let $k_0$ be its algebraic part over $\Q$. Let $v,w:k\hookrightarrow\C$ agree on $k_0$. Every finite diagram of smooth projective varieties over $k$ and algebraic correspondences has an integral Betti comparison between $v$ and $w$, compatible with cup products, K\"unneth maps, duality and those correspondences. If $k/k_0$ is finitely generated, one comparison exists simultaneously for all such varieties and correspondences. It extends to the induced rational tensor category and to the real scalar extensions of the integral realizations.
\end{result}
\begin{proof}
For a finite diagram, choose a finitely generated subfield $F\subset k$ over $k_0$ over which every equation, morphism and correspondence is defined. Because $k_0$ is algebraically closed in $k$, it is algebraically closed in $F$, and characteristic zero makes $F/k_0$ regular. A smooth model $S$ of $F/k_0$ is geometrically integral. After shrinking $S$, the varieties form smooth proper families, the correspondences extend, and all relations used in the finite diagram remain valid. The two embeddings define points of $S(\C)$ after the same embedding of $k_0$; they lie in every nonempty $k_0$-open subset because their maps on $F$ are injective. The complex manifold $S(\C)$ is connected and path connected. Parallel transport in the integral cohomology local systems along a path between those points gives the desired maps. Algebraic correspondences give flat maps of these local systems. Products and the orientation classes of complex fibres are transported, proving all stated compatibilities. Passing to the quotient by torsion and taking duals commute with these maps.

For the simultaneous assertion, take $F=k$ and one smooth model $S/k_0$. The field $k_0$ is countable. List all proper closed $k_0$-subvarieties of $S$; they form a countable family $Z_j$. Each variety and each correspondence under consideration extends over some open subset of $S$, so it is enough to join the two generic points by a path avoiding $\bigcup_j Z_j(\C)$. We give the needed topological argument. In a coordinate ball of complex dimension $d$, a proper complex analytic subset has real Hausdorff dimension at most $2d-2$. For a fixed point outside it, the set of endpoints whose joining segment meets it is contained in a countable union of cones of real Hausdorff dimension at most $2d-1$, and hence has $2d$-dimensional measure zero. This follows by exhausting the analytic subset by bounded pieces of its smooth strata and applying the locally Lipschitz map $(z,t)\mapsto x+t(z-x)$ on bounded intervals of $t$. Countable unions remain null. In each ball we can therefore choose an intermediate point whose two segments avoid all the $Z_j$. A finite chain of overlapping coordinate balls along any initial path, with overlap points chosen outside the same null sets, gives a path avoiding the entire union. Its single parallel-transport system proves the assertion. We use the smooth-proper local-system theorem in the form of Ehresmann's theorem, \cite[Theorem 12.1.1]{Arapura2003}; the argument just supplied is the simultaneous-choice step.
\end{proof}

The finite assertion concerns the specified finite diagram. A category called a ``finite tensor subcategory'' can contain further integral lattices and infinitely many divisibility conditions; one must specify these before identifying it with a finite diagram. The next counterexample is precisely sensitive to those conditions.

\begin{result}{E.1.2}{The universal countable integral assertion}{Refuted}
There is a countable algebraically closed field $k$ of transcendence degree one over $\Qb$, and two distinct complex places agreeing on $\Qb$, for which the integral Betti functors are not naturally isomorphic even on the abelian-variety and isogeny diagram. Each finite subdiagram admits an isomorphism. Thus the unrestricted inference on 15:221 from countability to a global integral isomorphism is false in this explicit interpretation.
\end{result}
\begin{proof}
Take $F=\Qb(t)$ and the elliptic curve
\[
 E/F:\quad y^2=x(x-1)(x-t).
\]
Fix a complex embedding $v$ with $v(t)=e$, and extend it to an algebraic closure $k=\overline F$. We identify the rational constants with the usual $\Qb\subset\C$. The field $k$ is countable, its algebraic part over $\Q$ is exactly $\Qb$, and $v$ is injective because $e$ is transcendental. The curve is smooth since $t(t-1)\ne0$.

Choose an integral homology basis for the Legendre family by transport from a point in $(0,1)$ to $e$. Its local monodromy around $t=0$, in an appropriate oriented basis, is
\[
 T_0=\begin{pmatrix}1&2\\0&1\end{pmatrix}.
\]
The exact integral matrix is recorded in \cite[Example 5.2, pp.~16--17]{OSU2024}. Riemann existence identifies the finite torsion monodromy with the Galois action on the corresponding torsion fields; we use \cite[Corollary 3.7]{OSU2024} and the algebraically closed base-field invariance in \cite[Theorem 1.1]{Landesman2024}, with characteristic zero and the smooth curve $\mathbf P^1-\{0,1,\infty\}$. In particular the image of
\[
 \Gal(F(E_{\mathrm{tors}})/F)\longrightarrow\SL_2(\widehat\Z)
\]
contains the closure of the powers of $T_0$, namely all matrices $u(2a)=\left(\begin{smallmatrix}1&2a\\0&1\end{smallmatrix}\right)$ for $a\in\widehat\Z$. This only uses those actual finite \emph{algebraic} torsion covers; no analytic cover is silently substituted for an algebraic one.

For an explicit noninteger $a$, define its residue modulo $n!$ by
\[
 a\bmod n!=s_n:=\sum_{m=1}^{n-1}m!\pmod {n!}.
\]
The residues are compatible, because $s_{n+1}-s_n=n!$. They define an element of $\widehat\Z$, since the factorial moduli are cofinal among positive integers. For $n\ge4$, $0<s_n<n!/2$, and $s_n\to\infty$. A fixed nonnegative integer has eventually constant least residue; a fixed negative integer has eventually least residue $n!-c>n!/2$. Neither agrees with $s_n$ for all $n$, so $a\notin\Z$. Choose the unique automorphism of the torsion field acting as $u(2a)$ and extend it to $\sigma\in\Gal(k/F)$. Put $w=v\circ\sigma$. These embeddings agree on $F$ but are distinct on the torsion field. They are not complex conjugate, since they both fix $i$.

Every finite cyclic subgroup $H\subset E[n]$ is defined over $k$. Form its algebraic quotient and quotient isogeny
\[
 f_H:E\longrightarrow E/H.
\]
All these varieties and morphisms belong to the concrete diagram. The complex curves $E_v$ and $E_w$ are identical because $v,w$ agree on $F$. Let $\Lambda=H_1(E_v,\Z)=H_1(E_w,\Z)$. Analytic uniformization identifies $f_H$ with the inclusion of lattices
\[
 \Lambda\subset\Lambda_H\subset n^{-1}\Lambda,
 \qquad \Lambda_H/\Lambda=H_v\subset n^{-1}\Lambda/\Lambda.
\]
For $w$ the corresponding subgroup is $H_w=u(2a)H_v$. These statements follow directly by identifying $E_v$ with $\C/\Lambda$ and taking the preimage of $H_v$ in $\C$; the quotient is $\C/\Lambda_H$ and its map is induced by the identity of $\C$.

A natural integral Betti isomorphism would have a component $A\in\GL_2(\Z)$ on $\Lambda$. Naturality for each $f_H$ forces
\[
 A\Lambda_H=\Lambda_{H_w}.
\]
Indeed the component on $E/H$ extends $A$ on the finite-index sublattice $\Lambda$, so its rational extension is forced to be $A$, and integral invertibility gives equality of the lattices. Modulo $n$, this says that $u(2a)^{-1}A$ stabilizes every cyclic subgroup of $(\Z/n)^2$. Stabilization of the three direct summands generated by $e_1,e_2,e_1+e_2$ forces a diagonal matrix with equal unit diagonal entries. Thus
\[
 u(2a)^{-1}A=\lambda I\quad\hbox{in }\GL_2(\widehat\Z),
 \qquad \lambda\in\widehat\Z^{\times}.
\]
Compatibility of the scalars modulo divisors follows from their uniqueness. Consequently $A_{21}=0$ and $A_{11}=A_{22}=\lambda$ are integers. Since $\det A=\pm1$, we obtain $\lambda=\pm1$. Now $A_{12}=\pm2a\in\Z$. An element of $\widehat\Z$ whose double is an integer is an integer: its rational value is integral in $\Z_2$, so that integer is even. Hence $a\in\Z$, the contradiction.

Every finite subdiagram descends to a finite extension of $F$, a finitely generated regular extension of $\Qb$. The two embeddings give points on its connected smooth complex model. E.1.1 supplies the required finite comparison. The finite matrix lemma and the factorial residues are checked by exact arithmetic in \texttt{code/verify\_places.py}; the proof above, rather than any finite bound in that program, establishes the infinite obstruction.
\end{proof}

\begin{result}{E.1.3}{The path-limit step and its completed reading}{Refuted; replacement proved}
Nonempty path sets in the tower on 15:221 need not have a nonempty inverse limit. The completed path torsors do admit compatible points, but those points act on profinite lattices; they need not preserve an integral lattice as E.1.2 demonstrates.
\end{result}
\begin{proof}
For an independent explicit path example set $k'=\bigcup_n\Qb(t^{1/n!})$, with $u_n=t^{1/n!}$ and $u_{n+1}^{n+1}=u_n$. Define
\[
 v(u_n)=\exp(1/n!),\qquad
 w(u_n)=\exp(1/n!+2\pi i s_n/n!).
\]
The residue compatibility makes these field embeddings fixing $\Qb$; their images of $t$ are both $e$. At level $n$, the power map $z\mapsto z^{n!}$ sends a path from $v(u_n)$ to $w(u_n)$ to a loop based at $e$ with winding number congruent to $s_n$ modulo $n!$. A compatible system would therefore have an integer winding number congruent to all $s_n$, impossible by the preceding proof. Each path set is nonempty. Transition maps on these path sets are not surjective: at the first level they miss the disallowed winding classes. The function-field transition has a nontrivial finite algebraic extension, so its generic fibre cannot be geometrically connected; shrinking models cannot change that field extension.

Replacing integer winding numbers by $\widehat\Z$ makes the compatible point $a$ available. More generally the path torsors of a profinite Galois category are nonempty compact spaces. Finite compatibility and compactness supply a point in their inverse limit without a countability requirement. Its scalar extension to $\Z_\ell$ gives the corresponding continuous lattice comparison for each prime. The map from integral transformations to these completed transformations is the completion map; E.1.2 proves that it need not be surjective. This covers both superscript readings $\pi^1$ and $\pi^{\wedge}$ without changing the intended coefficient ring silently.
\end{proof}

\begin{result}{E.1.4}{Uncountable fields: both behaviours occur}{Proved}
There are uncountable fields with distinct complex places inducing isomorphic integral Betti functors on all smooth projective varieties and algebraic correspondences. There are also uncountable fields with two places agreeing on $\Qb$ for which such an isomorphism does not exist.
\end{result}
\begin{proof}
For the positive example choose a transcendence basis $S\cup\{x\}$ of $\R/\Q$, with $S$ uncountable. Let $L$ be the real closure of $\Q(S)$ inside $\R$, put $K=L(i)$, and set $k=K(t)$. The field $K$ is algebraically closed, uncountable, and $x$ is transcendental over $K$. Define $v,w$ to fix $K$ and send $t$ to $x,x+1$, respectively. They are different complex places, both fixing $\Qb$. The polygonal path with vertices
\[
 x,\quad x+ix,\quad x+1+ix,\quad x+1
\]
avoids $K=L+iL$: along either vertical segment the real coordinate is outside $L$, and along the horizontal segment the imaginary coordinate is $x\notin L$. Every smooth projective variety or correspondence over $K(t)$ extends over some open subset of $\A^1_K$; its excluded complex points lie in $K$, since $K$ is algebraically closed. The one displayed path therefore gives parallel transport simultaneously for all of them. Cup products, correspondences and real scalar extensions are compatible as in E.1.1. Neither embedding has real image, so no extra real-place involution is being omitted.

For the negative example, enlarge the algebraically closed countable field $k$ of E.1.2 by an uncountable algebraically independent set and take an algebraic closure. Extend $v$ and $w$ through this extension, using the same independently chosen images for the new transcendence basis and then extending embeddings algebraically. Such choices exist inside $\C$ because its transcendence degree over a countable field is the continuum. Restriction of a natural integral comparison to the original abelian diagram would contradict E.1.2. The algebraic part remains $\Qb$. Thus cardinality alone gives no universal positive or negative assertion.
\end{proof}

\subsection{What the Hodge structure recovers}

\begin{result}{E.3.1}{The labelled integral elliptic test}{Proved}
Let $k$ be any characteristic-zero field admitting complex embeddings. On the diagram of elliptic curves over $k$, the integral Hodge realizations distinguish all \emph{chosen embeddings} $k\hookrightarrow\C$. For a number field, one labelled elliptic object suffices; for a finitely generated field, a finite labelled direct sum suffices. Complex places, defined modulo conjugation, are recovered after making the same quotient of the enriched data.
\end{result}
\begin{proof}
For $a\in k-\{0,1\}$ use the elliptic curve
\[
 E_a:\quad y^2=x^3-\frac{3a}{a-1}x-\frac{2a}{a-1}.
\]
Its discriminant is $1728a^2/(a-1)^3$ and its classical $j$-invariant is $1728a$, as checked symbolically in \texttt{verify\_places.py}. For an elliptic curve $E/\C$, its integral Hodge structure reconstructs its complex torus by the exact quotient
\[
 H^0(E,\Omega^1)^\vee\big/H_1(E,\Z),
 \qquad
 \delta\longmapsto\left(\omega\longmapsto\int_\delta\omega\right).
\]
The Hodge subspace supplies the dual vector space, and the integral pairing supplies the embedded lattice. An integral Hodge isomorphism therefore induces an isomorphism of the complex tori, hence of the elliptic curves, and preserves their $j$-invariants. For embeddings $v,w$, isomorphism of the two Hodge realizations of every $E_a$ forces $1728v(a)=1728w(a)$ for every $a\ne0,1$. They also agree on $0,1$, so $v=w$.

For a number field take a primitive element $a$, adjusting it by a rational integer so that $a\ne0,1$. Its single $E_a$ distinguishes all chosen embeddings. For $k=\Q(a_1,\ldots,a_r)$, apply the same adjustment to each generator and take the direct sum of their $H^1$ objects with the labelled inclusion and projection morphisms. A natural Hodge comparison commutes with each labelled projector and gives the individual isomorphisms; thus it forces equality on every generator. For a conjugate embedding the antiholomorphic comparison exchanges the Hodge labels and conjugates coefficients, exactly as in C.4.1. This identifies the quotient appropriate to unchosen complex places. No passage from integral isomorphism to rational isogeny is made in this proof.
\end{proof}

\begin{result}{E.3.2}{Failure of a universal single object in finite transcendence degree}{Proved}
For $k=\overline{\Qb(t)}$, no motive defined by finitely many algebraic equations, projectors and correspondences distinguishes all chosen complex embeddings through its realization alone. Nevertheless the entire labelled integral Hodge diagram distinguishes them by E.3.1.
\end{result}
\begin{proof}
Any one such object descends to a finite extension $F_0/\Qb(t)$. The algebraic closure $k/F_0$ has a nontrivial automorphism: the function field $F_0$ of a curve over $\Qb$ has, for example, a nonsquare function with a zero of odd order at a smooth point, giving a quadratic extension. Choose a nonidentity $\sigma\in\Gal(k/F_0)$ and set $w=v\circ\sigma$. They give distinct chosen embeddings and identical realizations of the descended object, with its defining finite diagram, because they agree on $F_0$. They fix $i$ and are therefore not related by complex conjugation. This proves the assertion for every individual object. The explicit nonsquare can be obtained from a local uniformizer at a smooth point by the approximation theorem on the curve; its odd valuation proves it cannot be a square.
\end{proof}

\paragraph{Remaining essential-image problem.}
E.3.1 proves the injectivity part of the Hodge-enriched reconstruction for every field in the prompt, including fields of finite transcendence degree that are not finitely generated. It does not identify all tensor functors equipped with unspecified ``structure suppl\'ementaire'' as geometric places. The Riemann--Hodge margin on 15:223 supplies no axioms for that target. Even after specifying integral polarizable Hodge realizations of all geometric objects and their labelled maps, the proof above recovers the values of an embedding from a functor already known to be geometric; it does not show that the values extracted from an arbitrary functor respect the field operations. This is a remaining exact obligation for E.2--E.3, not a surjectivity theorem.

\paragraph{Complex scalar extension and cardinality.}
For a specified neutral Tannakian category, Wibmer's theorem \cite[Theorems 1.1--1.2]{Wibmer2025} proves that two neutral fibre functors over an algebraically closed coefficient field are isomorphic, without a finite-type or countability restriction. Applying it after scalar extension to $\C$ gives the precise comparison map from integral isomorphism torsors to complex ones. The nonexistence in E.1.2 concerns a $\Z$-point; the complexified torsor can have a $\C$-point by this theorem. The theorem therefore closes the complex-coefficient cardinality question and leaves the explicit integral obstruction intact.
